% Chapter 4, Topic _Linear Algebra_ Jim Hefferon
%  http://joshua.smcvt.edu/linearalgebra
%  2001-Jun-12
\topic{Aturan Cramer}
\index{aturan Cramer|(}
% We have introduced determinant functions algebraically by looking
% for a formula to decide whether a matrix is nonsingular.
% After that introduction we saw a geometric interpretation, 
% that the determinant function
% gives the size of the box with sides formed by the columns of the matrix.
% Here we make a connection between the two views.

%<*CramersRuleExample0>
Sistem linear ekuivalen dengan relasi linear antarvektor.
\begin{equation*}
  \begin{linsys}{2}
     x_1  &+  &2x_2  &=  &6  \\
    3x_1  &+  &x_2   &=  &8 
  \end{linsys}
% \end{equation*}
% \begin{equation*}
  \qquad\Longleftrightarrow\qquad
  x_1\cdot\colvec{1 \\ 3}+x_2\cdot\colvec{2 \\ 1}=\colvec{6 \\ 8}
\end{equation*}
%</CramersRuleExample0>
Pada gambar di bawah, jajargenjang kecil dibentuk oleh vektor $\binom{1}{3}$ dan
$\binom{2}{1}$. Jajargenjang itu terletak di dalam jajargenjang dengan sisi
$x_1\binom{1}{3}$ dan~$x_2\binom{2}{1}$.
Menurut persamaan vektor, sudut terjauh jajargenjang yang lebih besar adalah
$\binom{6}{8}$.
\begin{center}
  \includegraphics{det/mp/ch4.1}
\end{center}
%<*CramersRuleExample1>
Gambar ini menyatakan ulang pertanyaan aljabar tentang pencarian solusi sistem
linear dalam istilah geometris:~dengan faktor $x_1$ dan~$x_2$ berapakah kita
harus mendilatasi sisi-sisi jajargenjang awal agar memenuhi jajargenjang lainnya?
%</CramersRuleExample1>

Kita dapat menggunakan gambar ini dan pemahaman geometris kita tentang
determinan untuk memperoleh rumus baru bagi penyelesaian sistem linear.
Bandingkan ukuran kotak-kotak berarsir berikut.
\begin{center}
   \includegraphics{det/mp/ch4.2}
   \hfil
   \includegraphics{det/mp/ch4.3}
   \hfil
   \includegraphics{det/mp/ch4.4}
\end{center}
%<*CramersRuleExample2>
Kotak kedua didefinisikan oleh vektor $x_1\binom{1}{3}$ dan $\binom{2}{1}$.
Salah satu sifat fungsi ukuran\Dash yaitu determinan\Dash menyatakan bahwa
ukuran kotak kedua adalah \( x_1 \) kali ukuran kotak pertama.
Kotak ketiga diperoleh dari kotak kedua melalui geseran, dengan menambahkan
$x_2\binom{2}{1}$ pada $x_1\binom{1}{3}$ untuk memperoleh
$x_1\binom{1}{3}+x_2\binom{2}{1}=\binom{6}{8}$,
bersama dengan $\binom{2}{1}$.
Determinan tidak dipengaruhi oleh geseran, sehingga ukuran kotak ketiga sama
dengan ukuran kotak kedua.
%</CramersRuleExample2>

%<*CramersRuleExample3>
Dengan menggabungkan semuanya, kita memperoleh
\begin{equation*}
  x_1\cdot \begin{vmat}[r]
     1  &2  \\
     3  &1
  \end{vmat}
  =
  \begin{vmat}
     x_1\cdot 1  &2  \\
     x_1\cdot 3  &1
  \end{vmat}
  =
  \begin{vmat}
     x_1\cdot 1+x_2\cdot 2  &2  \\
     x_1\cdot 3+x_2\cdot 1  &1
  \end{vmat}
  =
  \begin{vmat}[r]
     6  &2  \\
     8  &1
  \end{vmat}
\end{equation*}
%</CramersRuleExample3>
%<*CramersRuleExample4>
Penyelesaiannya memberikan nilai salah satu variabel.
\begin{equation*}
  x_1=
  \frac{\begin{vmat}[r]
     6  &2  \\
     8  &1
  \end{vmat} }{
  \begin{vmat}[r]
     1  &2  \\
     3  &1
  \end{vmat}  }
  =\frac{-10}{-5}=2
\end{equation*}
%</CramersRuleExample4>

Perumuman contoh ini adalah \definend{Aturan Cramer}:%
\index{determinan!aturan Cramer}%
\index{persamaan linear!solusi!aturan Cramer}
jika \( \deter{A}\neq 0 \), maka sistem \( A\vec{x}=\vec{b} \) memiliki solusi
tunggal
$
   x_i=\deter{B_i}/\deter{A}
$
dengan matriks $B_i$ dibentuk dari $A$ dengan mengganti kolom~$i$ oleh vektor
\( \vec{b} \). Pembuktiannya terdapat dalam
\nearbyexercise{ex:CramerRule}.

Sebagai contoh, untuk menyelesaikan sistem berikut bagi \( x_2 \),
\begin{equation*}
  \begin{mat}[r]
    1  &0  &4  \\
    2  &1  &-1 \\
    1  &0  &1
  \end{mat}
  \colvec{x_1 \\ x_2 \\ x_3}
  =\colvec[r]{2 \\ 1 \\ -1}
\end{equation*}
kita menjalankan perhitungan berikut.
\begin{equation*}
  x_2=
  \frac{ \begin{vmat}[r]
           1  &2  &4  \\
           2  &1  &-1 \\
           1  &-1 &1
         \end{vmat}  }{
         \begin{vmat}[r]
           1  &0  &4  \\
           2  &1  &-1 \\
           1  &0  &1
         \end{vmat}  }
  =\frac{-18}{-3}
\end{equation*}

Aturan Cramer memungkinkan kita menyelesaikan secara langsung sistem yang kecil
dan sederhana. Sebagai contoh, kita dapat menyelesaikan sistem dengan dua
persamaan dan dua variabel tak diketahui, atau tiga persamaan dan tiga variabel
tak diketahui, jika bilangan-bilangannya merupakan bilangan bulat kecil.
Kasus semacam itu cukup sering muncul sehingga banyak orang menganggap rumus ini
berguna.
 
Namun, penggunaan aturan ini untuk menyelesaikan sistem besar atau rumit tidak
praktis, baik dengan tangan maupun komputer.
Pendekatan yang didasarkan pada Metode Gauss lebih cepat.

\begin{exercises}
  \item 
    Gunakan Aturan Cramer untuk menyelesaikan masing-masing sistem bagi setiap
    variabelnya.
    \begin{exparts*}
      \partsitem $\begin{linsys}{2}
                    x  &- &y  &=  &4  \\
                   -x  &+ &2y &=  &-7
                  \end{linsys}$
      \partsitem $\begin{linsys}{2}
                    -2x  &+  &y  &=  &-2 \\
                      x  &-  &2y &=  &-2  
                  \end{linsys}$
    \end{exparts*}
    \begin{answer}
      \begin{exparts*}
        \partsitem Selesaikan setiap variabel secara terpisah.
          \begin{equation*}
            x=
             \frac{ \begin{vmat}[r]
                       4  &-1  \\
                      -7  &2   
                    \end{vmat}  }{
                    \begin{vmat}[r]
                       1  &-1  \\
                      -1  &2 
                    \end{vmat}  }
            =\frac{1}{1}=1
            \qquad
            y=
             \frac{ \begin{vmat}[r]
                       1  &4  \\
                      -1  &-7   
                    \end{vmat}  }{
                    \begin{vmat}[r]
                       1  &-1  \\
                      -1  &2 
                    \end{vmat}  }
            =\frac{-3}{1}=-3
          \end{equation*}
        \partsitem $x=2$, $y=2$
      \end{exparts*} 
    \end{answer}
  \item 
    Gunakan Aturan Cramer untuk menyelesaikan sistem ini bagi \( z \).
    \begin{equation*}
      \begin{linsys}{4}
        2x  &+  &y  &+  &z  &=  &1 \\
        3x  &   &   &+  &z  &=  &4 \\
         x  &-  &y  &-  &z  &=  &2
      \end{linsys}
    \end{equation*}
    \begin{answer}
      \( z=1 \)
    \end{answer}
  \item \label{ex:CramerRule}
    Buktikan Aturan Cramer.
    \begin{answer}
      Determinan tidak berubah oleh operasi kombinasi, termasuk kombinasi kolom,
      sehingga
      \( \det(B_i)=\det(\vec{a}_1,\dots,
      x_1\vec{a}_1+\dots+x_i\vec{a}_i+\dots+x_n\vec{a}_n,\dots,\vec{a}_n) \).
      Gunakan operasi mengambil $-x_1$ kali kolom pertama dan menambahkannya ke
      kolom ke-$i$, dan seterusnya, untuk melihat bahwa hasil ini sama dengan
      $\det(\vec{a}_1,\dots,x_i\vec{a}_i,\dots,\vec{a}_n)$.
      Selanjutnya, hasil tersebut sama dengan
      $x_i\cdot\det(\vec{a}_1,\dots,\vec{a}_i,\dots,\vec{a}_n)
         =x_i\cdot\det(A)$,
      seperti yang diperlukan.
    \end{answer}
  \item
   Berikut adalah pembuktian alternatif Aturan Cramer yang tidak secara eksplisit
   memuat geometri.
   Gunakan $X_i$ untuk menyatakan matriks identitas yang kolom~$i$-nya diganti
   oleh vektor variabel tak diketahui~$\vec{x}$, yaitu $x_1$, \ldots,~$x_n$.
   \begin{exparts}
     \partsitem Perhatikan bahwa $AX_i=B_i$.
     \partsitem Ambil determinan kedua ruas.
   \end{exparts}
   \begin{answer}
     \begin{exparts*}
       \partsitem
         Berikut adalah kasus sistem $\nbyn{2}$ dengan $i=2$.
         \begin{equation*}
           \begin{linsys}{2}
             a_{1,1}x_1 &+ &a_{1,2}x_2 &= &b_1 \\
             a_{2,1}x_1 &+ &a_{2,2}x_2 &= &b_2
           \end{linsys}
           \quad\Longleftrightarrow\quad
           \begin{mat}
             a_{1,1}  &a_{1,2}  \\
             a_{2,1}  &a_{2,2}
           \end{mat}
           \begin{mat}
             1  &x_1 \\
             0  &x_2
           \end{mat}
           =
           \begin{mat}
             a_{1,1}  &b_1 \\
             a_{2,1}  &b_2
           \end{mat}
         \end{equation*}
      \partsitem Fungsi determinan bersifat multiplikatif,
        $\det(B_i)=\det(AX_i)=\det(A)\cdot\det(X_i)$.
        Ekspansi Laplace menunjukkan bahwa $\det(X_i)=x_i$, dan penyelesaian
        bagi $x_i$ memberikan Aturan Cramer.
     \end{exparts*}
   \end{answer}

  \item 
    Misalkan suatu sistem linear memiliki persamaan sebanyak variabel tak
    diketahuinya, semua koefisien dan konstantanya merupakan bilangan bulat, dan
    matriks koefisiennya memiliki determinan~\( 1 \).
    Buktikan bahwa semua entri solusinya merupakan bilangan bulat.
    (\textit{Catatan.} Fakta ini sering digunakan untuk merancang sistem linear
    bagi latihan.)
    \begin{answer}
      Karena determinan $A$ tak nol, Aturan Cramer berlaku dan menunjukkan bahwa
      $x_i=\deter{B_i}/1$.
      Karena $B_i$ merupakan matriks bilangan bulat, determinannya juga bilangan
      bulat.
    \end{answer}
  \item 
    Gunakan Aturan Cramer untuk memberikan rumus solusi sistem linear dengan dua
    persamaan dan dua variabel tak diketahui.
    \begin{answer}
      Solusi bagi
      \begin{equation*}
        \begin{linsys}{2}
           ax  &+  by  &=  &e  \\
           cx  &+  dy  &=  &f  
        \end{linsys}
      \end{equation*}
      adalah
      \begin{equation*}
        x=\frac{ed-fb}{ad-bc}
        \qquad
        y=\frac{af-ec}{ad-bc}
      \end{equation*}
      dengan syarat, tentu saja, penyebutnya tidak nol.
    \end{answer}
  \item
    Dapatkah Aturan Cramer membedakan sistem tanpa solusi dari sistem dengan tak
    hingga banyak solusi?
    \begin{answer}
      Tidak selalu. Kedua jenis sistem bersifat singular, sehingga
      \( \deter{A}=0 \). Jika $\rank(A)=n-1$, sistem memiliki tak hingga banyak
      solusi tepat ketika semua \( \deter{B_i}=0 \); jika tidak, sekurang-kurangnya
      satu determinan tersebut tak nol dan sistem tidak konsisten.
      Namun, ketika $\rank(A)\leq n-2$, semua \( \deter{B_i} \) dapat bernilai
      nol meskipun sistem tidak konsisten. Sebagai contoh, $A=Z_{\nbyn{2}}$ dan
      $\vec{b}\neq\zero$ memberikan sistem tanpa solusi, tetapi
      $\deter{A}=\deter{B_1}=\deter{B_2}=0$.
    \end{answer}
  \item 
    Gambar pertama dalam Topik ini (yang tidak menggunakan determinan)
    menunjukkan kasus solusi tunggal.
    Buatlah gambar serupa bagi kasus tak hingga banyak solusi dan kasus tanpa
    solusi.
    \begin{answer}
      Kita dapat meninjau kedua kasus singular sekaligus melalui sistem
      \begin{equation*}
        \begin{linsys}{2}
           x_1  &+  &2x_2  &=  &6  \\
           x_1  &+  &2x_2  &=  &c 
        \end{linsys}
      \end{equation*}
      dengan $c=6$ tentu menghasilkan tak hingga banyak solusi, sedangkan setiap
      nilai $c$ lainnya menghasilkan sistem tanpa solusi.
      Persamaan vektor yang bersesuaian,
      \begin{equation*}
        x_1\cdot\colvec[r]{1 \\ 1}+x_2\cdot\colvec[r]{2 \\ 2}=\colvec[r]{6 \\ c}
      \end{equation*}
      memberikan gambar dua vektor yang bertumpang tindih.
      Keduanya terletak pada garis $y=x$.
      Dalam kasus $c=6$, vektor di ruas kanan juga terletak pada garis $y=x$,
      tetapi dalam kasus lainnya tidak.
    \end{answer}  
\end{exercises}
\index{aturan Cramer|)}
\endinput
